% Notation macros; independent of the document class (needs amsmath, amssymb, mathtools).
% Project notation is at the end (Min Cut-Path block, same as article 1).

% Number sets and expectation
\newcommand{\N}{\mathbb{N}}
\newcommand{\Z}{\mathbb{Z}}
\newcommand{\R}{\mathbb{R}}
\newcommand{\E}{\mathbb{E}}

% Operators: upright name, operator spacing
\newcommand{\OPT}{\mathrm{OPT}}
\DeclareMathOperator{\diam}{diam}
\DeclareMathOperator{\poly}{poly}

% Paired delimiters: \abs{x}; \abs*{x} scales to its content; \abs[\big]{x} sets the size by hand
\DeclarePairedDelimiter{\abs}{\lvert}{\rvert}
\DeclarePairedDelimiter{\ceil}{\lceil}{\rceil}
\DeclarePairedDelimiter{\floor}{\lfloor}{\rfloor}
\DeclarePairedDelimiter{\set}{\{}{\}}

% Names of computational problems in small caps: \textproblem{Min Cut-Path}.
% Not \problemname: llncs (and any class with a "problem" theorem environment) defines it.
\newcommand{\textproblem}[1]{\textsc{#1}}

% Min Cut-Path notation (\diam and \OPT are defined above)
\newcommand{\cp}{\operatorname{cp}}                  % cp(u,v): value of a minimum cut-path
\newcommand{\CP}{\operatorname{CP}}                  % CP(u,v): set of cut-paths
\newcommand{\MinCutPath}{\textsc{Min Cut-Path}}
\newcommand{\SSP}{\textsc{Separating Shortest Path}}
\newcommand{\ThreeSAT}{\textsc{3-SAT}}
